
\documentclass{article}
\usepackage{colortbl}
\usepackage{makecell}
\usepackage{multirow}
\usepackage{supertabular}

\begin{document}

\newcounter{utterance}

\centering \large Interaction Transcript for game `cladder', experiment `full\_v1.5\_default', episode 708 with Qwen3.6{-}27B.
\vspace{24pt}

{ \footnotesize  \setcounter{utterance}{1}
\setlength{\tabcolsep}{0pt}
\begin{supertabular}{c@{$\;$}|p{.15\linewidth}@{}p{.15\linewidth}p{.15\linewidth}p{.15\linewidth}p{.15\linewidth}p{.15\linewidth}}
    \# & \multicolumn{2}{c}{Player} && \multicolumn{2}{c}{Game Master} \\
    \hline

    \theutterance \stepcounter{utterance}  
    & & & \multicolumn{4}{p{0.6\linewidth}}{
        \cellcolor[rgb]{0.9,0.9,0.9}{
            \makecell[{{p{\linewidth}}}]{
                \texttt{\tiny{[P1$\langle$GM]}}
                \texttt{You are an expert in causal inference. The following question is not a typical commonsense query, but rather a meticulously designed question created by a professor specializing in causal inference, intended to assess the students' mastery of the course content.} \\
\texttt{Imagine a self{-}contained, hypothetical world with only the following conditions, and without any unmentioned factors or causal relationships: Gender has a direct effect on department competitiveness and admission status. Department competitiveness has a direct effect on admission status. The overall probability of male gender is 54\%. For individuals who are not male, the probability of admission acceptance is 49\%. For individuals who are male, the probability of admission acceptance is 24\%. Is admission acceptance less likely than admission rejection overall?} \\
\texttt{Start your answer with 'Yes' or 'No', followed by additional reasoning or evidence to support your explanation.} \\
            }
        }
    }
     \\ \\

    \theutterance \stepcounter{utterance}  
    & \multicolumn{4}{p{0.6\linewidth}}{
        \cellcolor[rgb]{0.561,0.737,0.561}{
            \makecell[{{p{\linewidth}}}]{
                \texttt{\tiny{[P1$\rangle$GM]}}
                \texttt{No} \\
\\ 
\texttt{To determine whether admission acceptance is less likely than admission rejection overall, we must calculate the total probability of admission acceptance, denoted as \$P(A)\$, using the Law of Total Probability.} \\
\\ 
\texttt{**1. Define the variables and given probabilities:**} \\
\texttt{*   Let \$G\$ represent Gender, where \$G=1\$ is Male and \$G=0\$ is Not Male (Female).} \\
\texttt{*   Let \$A\$ represent Admission Status, where \$A=1\$ is Acceptance and \$A=0\$ is Rejection.} \\
\\ 
\texttt{From the problem description, we have:} \\
\texttt{*   \$P(G=1) = 0.54\$ (Probability of being Male)} \\
\texttt{*   Therefore, \$P(G=0) = 1 {-} 0.54 = 0.46\$ (Probability of being Not Male)} \\
\texttt{*   \$P(A=1 | G=0) = 0.49\$ (Probability of acceptance given Not Male)} \\
\texttt{*   \$P(A=1 | G=1) = 0.24\$ (Probability of acceptance given Male)} \\
\\ 
\texttt{**2. Calculate the overall probability of admission acceptance \$P(A=1)\$:**} \\
\texttt{Using the Law of Total Probability:} \\
\texttt{\$\$P(A=1) = P(A=1 | G=1) \textbackslash{}cdot P(G=1) + P(A=1 | G=0) \textbackslash{}cdot P(G=0)\$\$} \\
            }
        }
    }
    & & \\ \\

    \theutterance \stepcounter{utterance}  
    & & & \multicolumn{2}{p{0.3\linewidth}}{
        \cellcolor[rgb]{0.9,0.9,0.9}{
            \makecell[{{p{\linewidth}}}]{
                \texttt{\tiny{[GM$|$GM]}}
                \texttt{No} \\
            }
        }
    }
    & & \\ \\

    \theutterance \stepcounter{utterance}  
    & & & \multicolumn{2}{p{0.3\linewidth}}{
        \cellcolor[rgb]{0.9,0.9,0.9}{
            \makecell[{{p{\linewidth}}}]{
                \texttt{\tiny{[GM$|$GM]}}
                \texttt{yes} \\
            }
        }
    }
    & & \\ \\

    \theutterance \stepcounter{utterance}  
    & & & \multicolumn{2}{p{0.3\linewidth}}{
        \cellcolor[rgb]{0.9,0.9,0.9}{
            \makecell[{{p{\linewidth}}}]{
                \texttt{\tiny{[GM$|$GM]}}
                \texttt{game\_result = LOSE} \\
            }
        }
    }
    & & \\ \\

\end{supertabular}
}

\end{document}
